\chapter*{Introducción}
\addcontentsline{toc}{chapter}{Introducción}

La seguridad ciudadana es una de las principales demandas de la población peruana y una responsabilidad de gestión de los gobiernos locales. En la región Cusco, los registros de la \gls{pnp} muestran un aumento de las denuncias por homicidio, y la \gls{enapres} sitúa al departamento entre los de mayor percepción de inseguridad del país \citep{inei2023enapres}. Sin embargo, los hechos no se distribuyen de manera uniforme: tienden a concentrarse en ciertos lugares y momentos \citep{ristea2020crimegis,medinaariza2023hotspots}, de modo que gestionarlos exige decidir dónde y cuándo intervenir, y no solo con qué recursos.

Para estudiar esa concentración se emplean mapas de calor, estimación de densidad kernel y técnicas de \termref{term-spatial-clustering}{agrupamiento espacial} \citep{ristea2020crimegis,hajela2020hotspot}. \gls{dbscan} \citep{choi2021mdstdbscan,xie2021significantdbscan} agrupa las ocurrencias según su densidad, no exige fijar el número de grupos y deja como ruido los puntos aislados; \gls{hdbscan} puede extender la idea a una jerarquía de densidades. Ambos dependen de parámetros cuya elección condiciona los grupos obtenidos, por lo que deben justificarse y verificarse. Además, que un algoritmo devuelva un grupo no demuestra que su concentración supere lo esperable por azar, y esa distinción es central en este plan.

El método ya se ha aplicado en el Perú \citep{cutid2025}, pero el extracto de ocurrencias del patrullaje municipal e integrado de San Sebastián se encuentra disponible como registros tabulares y carece de un flujo analítico local documentado que integre validación, agrupamiento, evaluación y presentación reproducible. Esta brecha es provisional y se contrasta con la revisión del capítulo ``Marco teórico y estado del arte''.

Por ello, el propósito de esta investigación es \parteProceso{caracterizar} los \parteContexto{incidentes de seguridad ciudadana registrados en SIPCOP Municipal en San Sebastián, Cusco}, mediante el desarrollo y evaluación de una \parteTecnologia{arquitectura de clústeres espaciales} que integre las etapas de validación y preparación de datos, agrupamiento, evaluación estadística, análisis temporal y presentación cartográfica. Se desarrollará una prueba de concepto reproducible y se aplicará al extracto disponible de \gls{sipcopm}, con 996 registros y fechas de alerta entre el 1 de enero y el 29 de agosto de 2025. El caso permitirá evaluar \gls{dbscan}, compararlo con \gls{hdbscan} y describir las concentraciones y sus distribuciones por turno, día de la semana y mes. El diagnóstico preliminar indica que las coordenadas tienen una resolución aproximada de 110 m (412 ubicaciones distintas en 996 registros), lo que condiciona el radio de agrupamiento; además, la tasa de registros cambia entre meses. La fecha de alerta no se asumirá como fecha exacta del hecho sin confirmar su definición. El producto de investigación será la arquitectura especificada y su prueba de concepto, no una plataforma municipal desplegada. Los resultados del caso se limitarán al extracto disponible y no presupondrán que existan concentraciones significativas.

El plan se organiza en cinco capítulos: 1, Generalidades; 2, Marco teórico y estado del arte; 3, Metodología; 4, Estudio de factibilidad; y 5, Aspectos administrativos. Al final se incluyen las referencias y los anexos.

\clearpage
